% Auto-extracted from background_photopolymerization.tex (body only)
\chapter{Chemical and process background}
\label{ch:background}

This chapter is the chemistry and process ontology of the campaign.
It is written for a reader with undergraduate physical chemistry and transport who has not necessarily worked on stereolithography.
Quantitative claims about a specific PEGDA formulation remain \stillreq{} until they are taken from a named experiment or paper figure under the same recipe.
Figures in this chapter are teaching cartoons, not experimental measurements.

\section{What we are trying to predict}

A photopolymerizable liquid is a mixture that stays liquid in the dark and becomes a covalently connected solid (a gel or a densely crosslinked plastic) where and when it is illuminated, provided enough of the right chemical events occur.
The engineering goal is to \emph{place} that solid: in a vat printer, as a layer of a part; in a microfluidic channel, as a fiber or a patch; in a hydrogel, as a spatial map of stiffness.

A tempting one-number description is the \emph{dose}
\begin{equation}
  E = \int I(t)\,\mathrm{d}t,
  \label{eq:dose}
\end{equation}
where $I$ is the irradiance (energy, or moles of photons, arriving per area per time) at a point on the resin.
Stereolithography practice has long used a threshold form of this idea: a working curve relating cure depth to $\ln(E/E_c)$ \citep{Jacobs1992book}.
That description is useful as a machine calibration for a \emph{given} resin, optics, and speed.
It is not a material state.
Two parcels of resin can receive the same $E$ and still differ in how many acrylate groups have reacted, whether a sample-spanning network exists, and how stiff the gel is.
The rest of this note explains why, by naming the chemical inventory, the competing rates, the transport of light and oxygen, and the difference between chemical conversion and a mechanically active network.

\paragraph{PEGDA as the running example.}
Poly(ethylene glycol) diacrylate (PEGDA) is a short chain with an acrylate group at each end.
Each acrylate carries a carbon--carbon double bond that can add to a radical.
Because there are two such groups per molecule, chain growth can connect different PEGDA molecules into a branched, then crosslinked, network \citep{Zhu2020,Montgomery2022,Slutzky2019}.
PEGDA is \emph{not} one resin: molecular weight (for example $\approx 250$ versus $\approx 575\,\mathrm{g\,mol^{-1}}$), water content, photoinitiator, and atmosphere change both the kinetics and the gel point.
Numbers fitted to one formulation must not be treated as universal PEGDA constants.

\section{The chemical inventory}

Figure~\ref{fig:species} is a cartoon of the species that matter for a typical free-radical acrylate photopolymerization.
Light does not ``glue PEGDA together'' in one step.
Light is absorbed by a \emph{photoinitiator}.
The initiator produces \emph{radicals}.
Radicals add to acrylate double bonds (\emph{propagation}), which is the step that consumes monomer and builds polymer.
Radicals also destroy one another (\emph{termination}) and are scavenged by dissolved oxygen (\emph{inhibition}).
Only the competition among these processes decides how much conversion occurs for a given illumination.

\begin{figure}[htbp]
  \centering
  \includegraphics[width=0.98\textwidth]{fig_background_species.pdf}
  \caption{Schematic chemical inventory for a free-radical PEGDA photopolymerization.
  Boxes are species classes, not a complete kinetic scheme.
  Conversion $p$ counts reacted acrylate groups; it is not yet gelation or modulus.}
  \label{fig:species}
\end{figure}

\subsection{Photoinitiators: Type~I versus Type~II}

Photoinitiators are classified by how they generate the first radicals after absorbing a photon \citep{Dadashi2016,Zhu2020}.

\begin{description}
  \item[Type~I]
    A single molecule undergoes bond cleavage and yields (typically two) initiating radicals.
    Irgacure~819 (phenylbis(2,4,6-trimethylbenzoyl)phosphine oxide), used by \citet{Montgomery2022} and in the illustrations of \citet{Dobson2024}, is in this class.
    A compact description of photolysis is
    \begin{equation}
      \mathrm{PI} + n\,h\nu \;\longrightarrow\; 2\,R^\bullet,
      \label{eq:typeI}
    \end{equation}
    with a rate proportional to the local irradiance and to the remaining initiator concentration \citep{Montgomery2022}.
  \item[Type~II]
    The absorber is excited to a triplet and must react with a co-initiator (often an amine) to produce radicals.
    Eosin~Y with triethanolamine, the system of \citet{Zhu2020}, is Type~II.
    The excited dye can also return to the ground state or be irreversibly bleached to a leuco product.
    Bleaching therefore changes \emph{both} the initiation rate \emph{and} the optical absorption of the sample as curing proceeds.
\end{description}

These are different chemical objects.
A rate law written for Irgacure cleavage is not interchangeable with an eosin~Y bleaching cascade, even if both are used to cure PEGDA.

\subsection{Monomer, conversion, and what FTIR actually measures}

Let $[M]$ be the molar concentration of unreacted acrylate \emph{functional groups} (not necessarily the concentration of PEGDA molecules).
The degree of conversion is
\begin{equation}
  p(t) = 1 - \frac{[M](t)}{[M]_0}.
  \label{eq:conversion}
\end{equation}
In practice $p$ is often inferred from the decay of the infrared acrylate band relative to an internal carbonyl standard \citep{Montgomery2022}.
That measurement reports \emph{how many double bonds have disappeared}.
It does not, by itself, report whether those additions connected two different chains or closed a loop on the same chain, nor whether the sample has gelled.

\subsection{Radicals as a short-lived catalyst, not a stored dose}

The radical concentration $[R^\bullet]$ is typically orders of magnitude smaller than $[M]$ or $[\mathrm{O}_2]$.
Radicals are generated by initiation and lost by termination and inhibition.
If generation and loss nearly balance, one may write a quasi-steady estimate such as
\begin{equation}
  [R^\bullet] \approx \sqrt{\frac{R_i}{2k_t}}
  \label{eq:qss}
\end{equation}
in the oxygen-free, bimolecular-termination limit, with $R_i$ the initiation rate and $k_t$ the termination coefficient \citep{Zhu2020}.
Equation~\eqref{eq:qss} is an \emph{approximation}, useful for insight.
It fails during the first instants after the lamp is switched on, when oxygen is still abundant, and whenever $k_t$ itself depends strongly on conversion.
Keeping the radical balance as a genuine time-dependent equation is therefore a modelling choice with chemical meaning: one refuses to assume that radicals forget their recent history \citep{Montgomery2022,Dobson2024,Slutzky2019}.

\subsection{Oxygen as a reagent, not a correction factor}

Molecular oxygen reacts rapidly with carbon-centred radicals to give comparatively unreactive peroxy species \citep{Decker1985,Dendukuri2008,Slutzky2019,Montgomery2022}:
\begin{equation}
  P^\bullet + \mathrm{O}_2 \;\longrightarrow\; P\mathrm{OO}^\bullet.
  \label{eq:oxygen}
\end{equation}
Until the local oxygen inventory is depleted, propagation cannot compete.
This has two engineering consequences.

First, there is an \emph{induction period}: illumination can proceed with almost no conversion while dissolved oxygen is consumed.
Second, oxygen is \emph{replenished} by diffusion from an interface (air, an oxygen-permeable window such as PDMS, or a flowing liquid that has not yet been illuminated).
\citet{Dendukuri2008} showed that in a PDMS microfluidic device the oxygen leaking through the elastomer sets a dead zone and a delay that cannot be read from a well-mixed bottle experiment.
\citet{Tumbleston2015} turned a persistent oxygen-inhibited liquid layer into a manufacturing feature (continuous liquid interface production).
The ontology point is the same: $\mathrm{O}_2$ is a transported chemical field with sources at boundaries, not a constant ``inhibition factor'' multiplying the dose.

\section{The kinetic skeleton (chain growth)}

A minimal reaction list, following the scheme used in flow by \citet{Slutzky2019} and in grayscale printing by \citet{Montgomery2022}, is:

\begin{center}
\begin{tabular}{@{}ll@{}}
\toprule
Photolysis / initiation & $\mathrm{PI} \xrightarrow{h\nu} R^\bullet$ \\
Chain initiation        & $R^\bullet + M \rightarrow P^\bullet$ \\
Propagation             & $P^\bullet_n + M \xrightarrow{k_p} P^\bullet_{n+1}$ \\
Bimolecular termination & $P^\bullet + P^\bullet \xrightarrow{k_t}$ dead polymer \\
Oxygen inhibition       & $P^\bullet + \mathrm{O}_2 \xrightarrow{k_O}$ inactive \\
\bottomrule
\end{tabular}
\end{center}

The polymerization rate (disappearance of acrylate) is then of the form
\begin{equation}
  R_p = k_p\,[R^\bullet]\,[M].
  \label{eq:Rp}
\end{equation}
Everything interesting is hidden in $[R^\bullet]$ and in the coefficients $k_p$ and $k_t$.

\subsection{\texorpdfstring{Why $k_p$ and $k_t$ are not constants}{Why kp and kt are not constants}}

As conversion rises, the mixture becomes more viscous.
Small monomers may still diffuse when large macroradicals can no longer find each other by centre-of-mass motion.
Termination therefore slows (\emph{autoacceleration}: the classical Trommsdorff--Norrish gel effect), while propagation later also becomes diffusion-limited (\emph{autodeceleration}).
This picture goes back at least to analyses of diffusion-controlled termination \citep{Benson1959} and is standard in thick-film photopolymerization models \citep{Goodner2002,Anseth1994}.

At still higher conversion, radicals can ``move'' by adding monomer rather than by translating: \emph{reaction-diffusion-controlled termination}, for which $k_t$ becomes proportional to $k_p[M]$ \citep{Anseth1994,Dobson2024,Montgomery2022}.
\citet{Dobson2024} emphasize a further refinement: short and long macroradicals do not share one $k_t$, because mobility depends on chain length and on whether the radical is already attached to the gel.

None of this is optional decoration if the process spans a wide conversion window.
It \emph{is} optional if one only asks whether a dilute flowing jet has consumed enough oxygen to gel at $\sim 2\%$ double-bond conversion \citep{Slutzky2019}.
The ontology must therefore allow \emph{reduced} kinetics for a narrowly posed question, without pretending that the reduced coefficients are the true molecular rate constants of PEGDA in every geometry.

\subsection{Thiol--ene is a different polymerization}

Acrylate chain growth is not the only photocurable chemistry.
Thiol--ene reactions have a step-growth character, different oxygen tolerance, and different network statistics \citep{Cramer2001}.
They must not be forced into the same four-species chain-growth skeleton without rewriting both the rates and the gelation criterion.
This note keeps them out of the PEGDA object.

\section{\texorpdfstring{Dose is not a state: the failure of exposure reciprocity}{Dose is not a state: the failure of exposure reciprocity}}

Photographic reciprocity says that the same $E=It$ produces the same blackening.
If photopolymerization were a one-step photoreaction with no competing second-order losses, something similar might hold for $p$.
Free-radical polymerization is not that reaction.

Bimolecular termination is second order in $[R^\bullet]$.
Higher irradiance produces a higher radical concentration, hence a higher \emph{fraction} of radicals lost to termination rather than to propagation.
Oxygen inhibition is also first order in $[R^\bullet]$.
The conversion therefore depends on the \emph{shape} of $I(t)$, not only on its integral.

\citet{Wydra2014} tested this on a BisGMA/TEGDMA dimethacrylate.
At equal dose, different irradiances produced different conversions and shrinkage stresses; the system ``does not and should not be expected'' to obey a reciprocity law, because bimolecular termination dominates.
Related kinetic analyses include the role of oxygen, viscosity, and time-varying intensity \citep{Lin2019}.
\citet{Dobson2024} show a complementary symptom: the apparent scaling of $R_p$ with $I$ can shift from roughly $\sqrt{I}$ (classic bimolecular termination) toward a more linear $I$ dependence when oxygen makes radical loss pseudo-first-order.

Figure~\ref{fig:reciprocity} is a \emph{cartoon} of that idea, not a plot of anyone's FTIR data.
The scientific statement is: if two illumination protocols are declared equivalent because $I_1 t_1 = I_2 t_2$, one has already assumed a constitutive law that acrylates typically violate.

\begin{figure}[htbp]
  \centering
  \includegraphics[width=0.82\textwidth]{fig_background_reciprocity.pdf}
  \caption{Schematic only.
  Two protocols with the same nominal dose need not yield the same conversion when termination is bimolecular and oxygen participates \citep{Wydra2014,Lin2019,Dobson2024}.
  The curves are not experimental measurements.}
  \label{fig:reciprocity}
\end{figure}

\section{Light as a field in a participating medium}

Irradiance is not a property of the lamp alone.
It is a field $I(\mathbf{x},t)$ inside an absorbing (and sometimes scattering) mixture.

\subsection{Beer--Lambert attenuation that evolves with chemistry}

If scattering and emission are negligible and the beam is treated as unidirectional, a steady radiative balance reduces to
\begin{equation}
  \frac{\partial I}{\partial z} + A(\mathbf{x},t)\,I = 0,
  \label{eq:beer}
\end{equation}
with $z$ along the beam \citep{Montgomery2022}.
The attenuation $A$ is a sum of contributions: remaining photoinitiator, unreacted monomer, polymer, and any dedicated absorber dye.
As initiator is consumed, $A$ can fall (\emph{photobleaching}, especially for Type~II dyes \citep{Zhu2020}).
As polymer forms, $A$ can rise if the polymer absorbs more than the monomer \citep{Montgomery2022}.
In an optically thick sample the result is a polymerization \emph{front} that advances from the illuminated face \citep{Cabral2005,Zhu2020}.

Equation~\eqref{eq:beer} is already an engineering reduction of radiative transfer.
It is appropriate for a clear PEGDA resin.
It is not appropriate for a ceramic-filled slurry, where scattering dominates.

\subsection{Pixels, grayscale, and the point-spread function}

A digital micromirror device (DMD) does not paint a perfectly sharp bitmap onto the resin.
Each micromirror illuminates a small neighbourhood.
\citet{Montgomery2022} represent that neighbourhood as a Gaussian and obtain the surface irradiance by summing all pixels, including grayscale (partial on-time of a mirror).
\citet{Meenakshisundaram2020} insist that the true kernel need not be Gaussian: measured irradiance from a projection head can have a flatter core and more bleed than a first-order Gaussian, because of mirror tilt, imperfect collimation, and the relay optics \citep[citing][]{Zhou2009}.

The ontology distinction:
\begin{itemize}
  \item the \emph{commanded bitmap} is a discrete pattern of grayscale values;
  \item the \emph{surface irradiance} is a continuous field obtained by convolution with the optical kernel and, if the head or the image moves, by a time integral;
  \item the \emph{in-resin irradiance} is that surface field after attenuation~\eqref{eq:beer};
  \item the \emph{chemical conversion} is the result of the kinetic skeleton driven by the in-resin irradiance.
\end{itemize}
Collapsing all four into a Jacobs working curve throws away the objects one needs in order to predict grayscale gradients and overcure \citep{Montgomery2022,Meenakshisundaram2020}.

\section{Two independent motions}

It is easy to say ``the process is moving.''
Two different velocities must be named (Fig.~\ref{fig:vel}).

\begin{figure}[htbp]
  \centering
  \includegraphics[width=0.98\textwidth]{fig_background_two_velocities.pdf}
  \caption{Left: the illumination pattern moves over a vat (scanning mask or scrolled image) while the resin is, to first approximation, still \citep{Meenakshisundaram2020}.
  Right: the resin flows through a fixed light window \citep{Slutzky2019}.
  A fluid element experiences the \emph{relative} motion of pattern and material.}
  \label{fig:vel}
\end{figure}

\paragraph{Pattern velocity $\mathbf{v}_{\mathrm{pattern}}$.}
In scanning-mask projection, the projector (or the bitmap on the DMD) translates so that a large area can be addressed with a small chip \citep{Meenakshisundaram2020}.
Frame rate, scan speed, and projected pixel pitch together decide how long a resin point sees each column of the movie.
One can even scroll the bitmap to \emph{cancel} stage motion, so that a feature is stationary in the vat while the hardware moves.

\paragraph{Material velocity $\mathbf{u}$.}
In a microfluidic jet, oligomer is carried through a stationary UV spot of length $L$ at speed $U$ \citep{Slutzky2019}.
The residence time in the light is $t_L = L/U$.
If the lamp is pulsed for a duration $t_{\mathrm{uv}}$, the relevant dimensionless exposure is $t_{\mathrm{uv}}/t_L$.

A material element at position $\mathbf{x}_p(t)$ with $\mathrm{d}\mathbf{x}_p/\mathrm{d}t=\mathbf{u}$ sees
\begin{equation}
  I_p(t) = I\bigl(\mathbf{x}_p(t),t\bigr).
\end{equation}
That sampled history is what the kinetic skeleton should be fed.
Calling both ``speed'' hides whether one is moving the lamp, the mask, or the liquid.

\section{Flow replenishes oxygenated resin}

In a closed, stagnant film, dissolved oxygen is a finite inventory (plus whatever diffuses from the walls).
In a continuous flow, fresh, oxygenated liquid continually enters the illumination zone.
\citet{Slutzky2019} make this the condition for fiber formation: polymerization remains negligible until oxygen is depleted, and the critical processing speed scales with how fast photolysis can eat the incoming $\mathrm{O}_2$ inventory,
\begin{equation}
  U_c \sim \frac{L\,\phi\,\varepsilon I\,[\mathrm{PI}]_0}{[\mathrm{O}_2]_0},
  \label{eq:Uc}
\end{equation}
up to the precise photophysical factors in their analysis.
They also argue that, for their jet diameter and speeds, oxygen \emph{diffusion into the jet during the residence time} is a small correction relative to advection of the dissolved inventory.
That statement is geometry-specific.
In a thin film against PDMS, the opposite limit can hold: diffusion from the wall dominates \citep{Dendukuri2008}.

The ontology: \emph{advection} and \emph{diffusion} of oxygen are different supply routes.
A model that includes only one of them is answering a different experiment.

\section{Conversion is not gelation is not modulus}

Figure~\ref{fig:stack} is the single most important hierarchy in this note.

\begin{figure}[htbp]
  \centering
  \includegraphics[width=0.88\textwidth]{fig_background_network_stack.pdf}
  \caption{Four layers that must not be identified.
  Light history and species fields determine conversion $p$; $p$ and composition determine whether a sample-spanning network exists and how many chains are elastically active; modulus is a mechanical consequence of that network \citep{Zhu2020,Macosko1976,Zhong2016}.}
  \label{fig:stack}
\end{figure}

\subsection{Gelation}

Gelation is the appearance of a sample-spanning covalently connected cluster: the weight-average molecular weight diverges \citep{Flory1941,Macosko1976}.
For a chain-growth polymerization of a multifunctional acrylate, that can occur at \emph{low} conversion because each chain can connect many monomers.

Two operational definitions in the present corpus must not be mixed without comment.

\begin{itemize}
  \item \citet{Slutzky2019} take gelation, following earlier microfluidic work, when about $2\%$ of double bonds have reacted ($[M]/[M]_0=0.98$).
    That is a convenient kinetic threshold for ``the jet no longer behaves as a liquid'' under their conditions.
  \item \citet{Zhu2020} compute a conversion at gelation from a recursive Macosko--Miller argument that includes intramolecular loops.
    If $\theta$ is the probability that a reacted acrylate connected \emph{inter}molecularly rather than as a loop, then
    \begin{equation}
      P_{\mathrm{gel}} = \frac{1-\theta}{2\theta},
      \qquad
      \theta = \frac{[\mathrm{PEGDA}]_0}{[\mathrm{PEGDA}]_0 + 2C_0},
      \label{eq:pgel}
    \end{equation}
    so $P_{\mathrm{gel}} = C_0/[\mathrm{PEGDA}]_0$.
    If the precursor is more dilute than $C_0$, $P_{\mathrm{gel}}\ge 1$: the sample never gels even at full acrylate conversion, because too many additions are wasted as loops.
\end{itemize}

Both can be legitimate.
They are not the same predicate.

\subsection{Loops and elastically active density}

Even after gelation, not every reacted acrylate contributes a load-bearing chain.
An intramolecular loop links a strand to itself and does not transmit stress the way an intermolecular crosslink does \citep{Zhong2016,Zhu2020}.
\citet{Zhu2020} construct a fractional crosslinking density $\eta$ from the probabilities that arms lead to the infinite network, and then write a phantom-network shear modulus
\begin{equation}
  G = \eta\,N_A\,[\mathrm{PEGDA}]_0\,k_B T.
  \label{eq:G}
\end{equation}
The ideal affine or phantom network without loops is an upper bound; real PEGDA gels sit below it.
Fitting $G$ as a smooth function of dose, skipping $\eta$, hides the concentration-dependent gel point that \citet{Zhu2020} used to design stiffness gradients.

\subsection{Heat, shrinkage, and swelling}

Polymerization is exothermic; thick or intensely illuminated samples heat, and rates depend on temperature \citep{Goodner2002,Dobson2024}.
Conversion also shrinks the volume and, in hydrogels, the network later imbibes solvent.
Those are real couplings.
They are \emph{downstream} of a trustworthy conversion field.
Treating them as the first objects to predict, before $p(\mathbf{x},t)$ is even defined, inverts the hierarchy of Fig.~\ref{fig:stack}.

\section{How the five campaign papers partition the ontology}

The local literature folder is not five versions of one model.
Each paper owns a different layer.

\begin{table}[htbp]
\centering
\caption{Division of labour among the five unique reports in the campaign corpus.
``Owns'' means: this is the layer the paper actually closes with equations and experiments.
It does not mean the other layers are false.}
\label{tab:papers}
\small
\begin{tabular}{@{}p{0.28\textwidth}p{0.62\textwidth}@{}}
\toprule
Report & Layer owned \\
\midrule
\citet{Dobson2024} &
High-resolution kinetic ontology: many species, free-volume transport, short versus long radicals, oxygen, heat, one- and three-dimensional balances.
Chemistry illustrated on HDDA, not PEGDA.
\\
\citet{Montgomery2022} &
Spatial grayscale printing: four tracked species, conversion-dependent $k_p$ and $k_t$, evolving Beer--Lambert field, pixel overlap, lateral diffusion that blurs commanded edges.
PEGDA~$250$ / Irgacure~819.
\\
\citet{Meenakshisundaram2020} &
Projection kinematics: measured (possibly non-Gaussian) pixel irradiance, scanning mask, discrete frames, dose accumulation.
Cure profile compared to printed dimensions; chemistry is not the species skeleton above.
\\
\citet{Slutzky2019} &
Flowing material through a fixed light window: advection of initiator, radicals, monomer, and oxygen; critical speed; pulse duration versus residence time.
PEGDA~$575$ jet; oxygen diffusion into the jet argued small.
\\
\citet{Zhu2020} &
Type~II eosin~Y bleaching plus conversion-to-gelation-to-modulus closure with loops.
Frontal photopolymerization in depth.
PEGDA~$575$ hydrogel; no advection.
\\
\bottomrule
\end{tabular}
\end{table}

A predictive description of a \emph{flowing grayscale PEGDA hydrogel} needs all four chemical--physical layers---kinetics, optics, transport, network---but it should inherit each layer from the paper that actually resolved it, not average their rate constants.

\section{What must not be identified}

For a masters-level reader, the practical discipline is a list of forbidden equivalences.

\begin{enumerate}
  \item $E=\int I\,\mathrm{d}t$ $\not\equiv$ conversion $p$. Reciprocity fails when termination is bimolecular and oxygen participates \citep{Wydra2014,Lin2019}.
  \item Commanded bitmap $\not\equiv$ surface irradiance $\not\equiv$ in-resin $I(\mathbf{x},t)$ \citep{Meenakshisundaram2020,Montgomery2022}.
  \item Pattern velocity $\not\equiv$ resin velocity \citep{Meenakshisundaram2020,Slutzky2019}.
  \item Dissolved oxygen inventory $\not\equiv$ oxygen leaking through a wall \citep{Slutzky2019,Dendukuri2008}.
  \item Type~I photolysis $\not\equiv$ Type~II dye/amine bleaching \citep{Dadashi2016,Zhu2020,Montgomery2022}.
  \item PEGDA~$250$ in a dry thin film $\not\equiv$ PEGDA~$575$ in water $\not\equiv$ HDDA bulk \citep{Montgomery2022,Zhu2020,Dobson2024}.
  \item Acrylate conversion $p$ $\not\equiv$ gelation $\not\equiv$ elastically active density $\eta$ $\not\equiv$ modulus $G$ \citep{Zhu2020,Flory1941,Zhong2016}.
  \item A $2\%$ double-bond threshold \citep{Slutzky2019} $\not\equiv$ the Macosko--Miller $P_{\mathrm{gel}}([\mathrm{PEGDA}]_0)$ \citep{Zhu2020}.
  \item Chain-growth acrylate kinetics $\not\equiv$ thiol--ene step growth \citep{Cramer2001}.
  \item Jacobs working-curve depth \citep{Jacobs1992book} $\not\equiv$ a species-resolved polymer state.
\end{enumerate}

\section{Implications for a predictive description}

The ontology above already implies an order of questions, still without reference to implementation.

\begin{enumerate}
  \item \textbf{Local constitutive response.}
    For a prescribed $I(t)$, how do initiator, radicals, oxygen, and acrylate evolve?
    This is the question \citet{Dobson2024} answer at high chemical resolution and \citet{Montgomery2022} answer with four species and conversion-dependent coefficients.
    Until this map is trusted---including an equal-dose test that is \emph{allowed to fail}---spatial complexity only conceals error.
  \item \textbf{Optical boundary condition.}
    What irradiance actually sits on $z=0$ as a function of bitmap, pixel kernel, and pattern motion \citep{Meenakshisundaram2020}?
    Then, how does $A(p,[\mathrm{PI}])$ carry that field into the depth \citep{Montgomery2022,Zhu2020}?
  \item \textbf{Transport.}
    Which species move relative to the material?
    Oxygen almost always does; monomer may; network-bound radicals do not \citep{Dobson2024,Montgomery2022}.
    Is the supply advective \citep{Slutzky2019} or diffusive from a wall \citep{Dendukuri2008}?
  \item \textbf{Network closure.}
    Given $p$ and composition, has the sample gelled, and what is $\eta$ and $G$ \citep{Zhu2020}?
    This step must not be used to patch kinetic coefficients.
\end{enumerate}

The campaign choice that follows from this background is conservative: treat PEGDA chain-growth as the first chemical object; treat Dobson's full species set as a \emph{reference} against which reduced $k_p(p)$ and $k_t(p)$ can be judged, not as the first process model; treat scanning projection and plug flow as two different kinematic experiments; treat Zhu's loop-aware modulus as a post-processor on conversion.
Those are scientific restrictions, not software preferences.

